\documentclass[11pt,a4paper]{article}

\usepackage[margin=1in]{geometry}
\usepackage{amsmath}
\usepackage{amssymb}
\usepackage{bm}
\usepackage{enumitem}

\title{VNAV LAB2 -- ROS and Rigid Transformations}
\author{Junxi Li}
\date{}

\begin{document}

\maketitle

\section{Deliverable 1: Nodes, Topics, and Launch Files}

\subsection{Nodes in the Static Scenario}

Ignoring the \texttt{/rosout} node, the nodes running in the static
two-drone scenario are

\begin{verbatim}
/av1broadcaster
/av2broadcaster
/plots_publisher_node
/rviz
\end{verbatim}

The two broadcaster nodes publish the static transforms from the world
frame to the two aerial-vehicle frames. The plots publisher generates
the drone meshes and trajectory markers, while RViz visualizes the
published information.

\subsection{Running the Static Scenario without roslaunch}

The same static scenario can be launched manually using separate
terminals.

Terminal 1:
\begin{verbatim}
roscore
\end{verbatim}

Terminal 2:
\begin{verbatim}
rosrun tf2_ros static_transform_publisher \
1 0 0 0 0 0 1 world av1 __name:=av1broadcaster
\end{verbatim}

Terminal 3:
\begin{verbatim}
rosrun tf2_ros static_transform_publisher \
0 0 1 0 0 0 1 world av2 __name:=av2broadcaster
\end{verbatim}

Terminal 4:
\begin{verbatim}
rosrun two_drones_pkg plots_publisher_node
\end{verbatim}

Terminal 5:
\begin{verbatim}
rviz -d $(rospack find two_drones_pkg)/config/default.rviz
\end{verbatim}

\subsection{ROS Topics}

The node \texttt{/av1broadcaster} publishes the transform between
\texttt{world} and \texttt{av1} on \texttt{/tf\_static}. Similarly,
\texttt{/av2broadcaster} publishes the transform between
\texttt{world} and \texttt{av2}.

The node \texttt{/plots\_publisher\_node} obtains transform information
from TF and publishes a message of type
\texttt{visualization\_msgs/MarkerArray} on the topic
\texttt{/visuals}.

RViz subscribes to the TF information and to \texttt{/visuals}. The
topic that causes the two drone meshes to appear in RViz is therefore

\[
\boxed{\texttt{/visuals}}.
\]

\subsection{Effect of Omitting static:=True}

The launch file defines the argument \texttt{static} with default value
\texttt{false}. When \texttt{static:=True}, the nodes
\texttt{/av1broadcaster} and \texttt{/av2broadcaster} are launched and
publish fixed transforms.

When the argument is omitted, its value is false. The launch-file
\texttt{if} condition therefore disables the two static broadcasters,
while the \texttt{unless} condition enables
\texttt{frames\_publisher\_node}. Consequently, the transforms become
time-varying and the two drones move according to the trajectories
defined in the problem.

\section{Deliverable 4: Mathematical Derivations}

Let

\[
c = \cos t,
\qquad
s = \sin t.
\]

The positions of AV1 and AV2 in the world frame are

\[
{}^{w}\bm{o}_1 =
\begin{bmatrix}
c\\
s\\
0
\end{bmatrix},
\qquad
{}^{w}\bm{o}_2 =
\begin{bmatrix}
s\\
0\\
\cos 2t
\end{bmatrix}.
\]

AV1 has zero roll and pitch and yaw equal to $t$. Therefore,

\[
{}^{w}\bm{R}_1 =
\begin{bmatrix}
c & -s & 0\\
s & c & 0\\
0 & 0 & 1
\end{bmatrix}.
\]

AV2 has the same orientation as the world frame.

\subsection{AV2 Trajectory in the World Frame}

The coordinates of AV2 satisfy

\[
x_2^w = \sin t,
\qquad
y_2^w = 0,
\qquad
z_2^w = \cos 2t.
\]

Using the identity

\[
\cos 2t = 1 - 2\sin^2 t,
\]

we obtain

\[
z_2^w = 1 - 2(x_2^w)^2.
\]

Hence

\[
\boxed{
z_2^w = 1 - 2(x_2^w)^2,
\qquad
y_2^w = 0
}
\]

which is a parabola in the $x_w$-$z_w$ plane.

\subsection{Position of AV2 Relative to AV1}

The homogeneous transformation of AV1 with respect to the world frame is

\[
{}^{w}\bm{T}_1 =
\begin{bmatrix}
c & -s & 0 & c\\
s & c & 0 & s\\
0 & 0 & 1 & 0\\
0 & 0 & 0 & 1
\end{bmatrix}.
\]

Since AV2 does not rotate with respect to the world frame,

\[
{}^{w}\bm{T}_2 =
\begin{bmatrix}
1 & 0 & 0 & s\\
0 & 1 & 0 & 0\\
0 & 0 & 1 & \cos 2t\\
0 & 0 & 0 & 1
\end{bmatrix}.
\]

The relative transformation is

\[
{}^{1}\bm{T}_2 =
({}^{w}\bm{T}_1)^{-1}
{}^{w}\bm{T}_2.
\]

For the translational component,

\[
{}^{1}\bm{o}_2
=
({}^{w}\bm{R}_1)^T
\left(
{}^{w}\bm{o}_2 -
{}^{w}\bm{o}_1
\right).
\]

First,

\[
{}^{w}\bm{o}_2 -
{}^{w}\bm{o}_1
=
\begin{bmatrix}
s-c\\
-s\\
\cos 2t
\end{bmatrix}.
\]

Furthermore,

\[
({}^{w}\bm{R}_1)^T =
\begin{bmatrix}
c & s & 0\\
-s & c & 0\\
0 & 0 & 1
\end{bmatrix}.
\]

Therefore,

\[
{}^{1}\bm{o}_2 =
\begin{bmatrix}
c & s & 0\\
-s & c & 0\\
0 & 0 & 1
\end{bmatrix}
\begin{bmatrix}
s-c\\
-s\\
\cos 2t
\end{bmatrix}.
\]

After multiplication,

\[
\boxed{
{}^{1}\bm{o}_2(t) =
\begin{bmatrix}
\sin t\cos t - 1\\
-\sin^2 t\\
\cos 2t
\end{bmatrix}
}.
\]

\subsection{Plane Containing the Relative Trajectory}

From the previous result,

\[
y_2^1 = -\sin^2 t
\]

and

\[
z_2^1 = \cos 2t
       = 1 - 2\sin^2 t.
\]

Thus

\[
z_2^1 = 1 + 2y_2^1.
\]

Hence every point of the relative trajectory satisfies

\[
\boxed{
z_2^1 - 2y_2^1 - 1 = 0
}.
\]

The relative trajectory therefore lies entirely in the plane

\[
\boxed{
\Pi:\; z - 2y - 1 = 0
}.
\]

\subsection{A Two-Dimensional Coordinate Frame on the Plane}

Following the suggested centre, define

\[
{}^{1}\bm{p} =
\begin{bmatrix}
-1\\
-\frac{1}{2}\\
0
\end{bmatrix}.
\]

The displacement from this point to AV2 is

\[
\bm{q}
=
{}^{1}\bm{o}_2 -
{}^{1}\bm{p}.
\]

Substitution gives

\[
\bm{q}
=
\begin{bmatrix}
\sin t\cos t\\
\frac{1}{2}-\sin^2 t\\
\cos 2t
\end{bmatrix}.
\]

Using the double-angle identities,

\[
\sin t\cos t = \frac{1}{2}\sin 2t
\]

and

\[
\frac{1}{2}-\sin^2 t
=
\frac{1}{2}\cos 2t,
\]

so

\[
\bm{q}
=
\begin{bmatrix}
\frac{1}{2}\sin 2t\\
\frac{1}{2}\cos 2t\\
\cos 2t
\end{bmatrix}.
\]

Choose the orthonormal axes

\[
\bm{e}_{x_p}
=
\begin{bmatrix}
1\\0\\0
\end{bmatrix},
\]

\[
\bm{e}_{y_p}
=
\frac{1}{\sqrt{5}}
\begin{bmatrix}
0\\1\\2
\end{bmatrix},
\]

and

\[
\bm{e}_{z_p}
=
\frac{1}{\sqrt{5}}
\begin{bmatrix}
0\\-2\\1
\end{bmatrix}.
\]

The rotation matrix from frame $p$ to frame 1 is therefore

\[
{}^{1}\bm{R}_{p}
=
\begin{bmatrix}
1 & 0 & 0\\
0 & \frac{1}{\sqrt{5}} & -\frac{2}{\sqrt{5}}\\
0 & \frac{2}{\sqrt{5}} & \frac{1}{\sqrt{5}}
\end{bmatrix}.
\]

The coordinates in the new frame are

\[
{}^{p}\bm{o}_2
=
({}^{1}\bm{R}_{p})^T
\bm{q}.
\]

This yields

\[
\boxed{
{}^{p}\bm{o}_2(t)
=
\begin{bmatrix}
\frac{1}{2}\sin 2t\\
\frac{\sqrt{5}}{2}\cos 2t\\
0
\end{bmatrix}
}.
\]

The zero third component confirms that the trajectory lies in the
$x_p$-$y_p$ plane.

\subsection{Ellipse and Semi-Axes}

From

\[
x_p = \frac{1}{2}\sin 2t
\]

and

\[
y_p = \frac{\sqrt{5}}{2}\cos 2t,
\]

we obtain

\[
\frac{x_p^2}{(1/2)^2}
=
\sin^2 2t
\]

and

\[
\frac{y_p^2}{(\sqrt{5}/2)^2}
=
\cos^2 2t.
\]

Adding both expressions,

\[
\boxed{
\frac{x_p^2}{1/4}
+
\frac{y_p^2}{5/4}
=
1
}.
\]

Therefore, the relative trajectory is an ellipse centred at the origin
of the $p$ frame.

Its semi-axis lengths are

\[
\boxed{
a = \frac{\sqrt{5}}{2},
\qquad
b = \frac{1}{2}
}.
\]

\section{Deliverable 5: More Properties of Quaternions}

Let a quaternion be represented as

\[
q =
\begin{bmatrix}
q_1&q_2&q_3&q_4
\end{bmatrix}^T,
\]

where the first three components form the vector part and $q_4$ is the
scalar part.

Let

\[
e =
\begin{bmatrix}
0\\0\\0\\1
\end{bmatrix}
\]

denote the unit quaternion representing the identity rotation.

The quaternion product satisfies

\[
q_a \otimes q_b
=
\Omega_1(q_a)q_b
=
\Omega_2(q_b)q_a.
\]

\subsection{Orthogonality of $\Omega_1$ and $\Omega_2$}

Quaternion norms are multiplicative:

\[
\|q_a\otimes q_b\|
=
\|q_a\|\|q_b\|.
\]

For a unit quaternion $q$ and arbitrary $x\in\mathbb{R}^4$,

\[
\|\Omega_1(q)x\|
=
\|q\otimes x\|
=
\|q\|\|x\|
=
\|x\|.
\]

Thus $\Omega_1(q)$ preserves the Euclidean norm of every vector.
Therefore,

\[
x^T\Omega_1(q)^T\Omega_1(q)x
=
x^Tx
\]

for every $x$, which implies

\[
\boxed{
\Omega_1(q)^T\Omega_1(q)=I_4
}.
\]

Since $\Omega_1(q)$ is square,

\[
\boxed{
\Omega_1(q)\Omega_1(q)^T=I_4
}.
\]

Similarly,

\[
\|\Omega_2(q)x\|
=
\|x\otimes q\|
=
\|x\|\|q\|
=
\|x\|,
\]

so

\[
\boxed{
\Omega_2(q)^T\Omega_2(q)
=
\Omega_2(q)\Omega_2(q)^T
=
I_4
}.
\]

Intuitively, multiplication by a unit quaternion preserves quaternion
norm. Hence the corresponding linear transformations in
$\mathbb{R}^4$ must be orthogonal.

\subsection{Mapping $q$ to the Identity Quaternion}

Since the identity quaternion satisfies

\[
q\otimes e=q,
\]

we have

\[
q=\Omega_1(q)e.
\]

Because $\Omega_1(q)$ is orthogonal,

\[
\Omega_1(q)^{-1}
=
\Omega_1(q)^T.
\]

Therefore,

\[
\boxed{
\Omega_1(q)^Tq=e
}.
\]

Similarly,

\[
e\otimes q=q
\]

implies

\[
q=\Omega_2(q)e.
\]

Hence,

\[
\boxed{
\Omega_2(q)^Tq=e
}.
\]

Consequently,

\[
\boxed{
\Omega_1(q)^Tq
=
\Omega_2(q)^Tq
=
\begin{bmatrix}
0\\0\\0\\1
\end{bmatrix}
}.
\]

\subsection{Commutativity of the Two Linear Operators}

Let $x,y,z\in\mathbb{R}^4$. Then

\[
\Omega_1(x)\Omega_2(y)z
=
x\otimes(z\otimes y).
\]

Quaternion multiplication is associative, so

\[
x\otimes(z\otimes y)
=
(x\otimes z)\otimes y.
\]

Therefore,

\[
\Omega_1(x)\Omega_2(y)z
=
\Omega_2(y)\Omega_1(x)z.
\]

Since this holds for every $z$,

\[
\boxed{
\Omega_1(x)\Omega_2(y)
=
\Omega_2(y)\Omega_1(x)
}.
\]

Now define the quaternion conjugate

\[
\bar{y}
=
\begin{bmatrix}
-y_1\\
-y_2\\
-y_3\\
y_4
\end{bmatrix}.
\]

From the definition of $\Omega_2$,

\[
\Omega_2(y)^T
=
\Omega_2(\bar{y}).
\]

Applying the first commutation result with $\bar{y}$ gives

\[
\Omega_1(x)\Omega_2(\bar{y})
=
\Omega_2(\bar{y})\Omega_1(x).
\]

Therefore,

\[
\boxed{
\Omega_1(x)\Omega_2(y)^T
=
\Omega_2(y)^T\Omega_1(x)
}.
\]

\section{Conclusion}

The ROS experiments verified the expected rigid-body transformations in
both the world frame and the moving AV1 frame. The TF-based
visualization produced a circular trajectory for AV1, a parabolic
trajectory for AV2 in the world frame, and an elliptical relative
trajectory for AV2 when expressed in AV1's frame. The analytical
derivations agree with the observed RViz results.

\end{document}
